% A l'original hi ha una «X» solta al costat de l'etiqueta «a)»; no s'ha transcrit.
\begin{solapartat}
\figura[0.6]{sol_fig1.png}

\punts{0,1} Esquema, $\cos\theta = \dfrac{2,50}{3,00}$

\punts{0,2}
\begin{align*}
E &= k\frac{Q}{r^2}\\
E_x &= E_1\cos\theta + E_2\cos\theta\\
E_y &= E_1\sin\theta - E_2\sin\theta
\end{align*}

\punts{0,2}
\begin{align*}
E_x &= \frac{2,50}{3,00}\times\frac{8,99\times 10^{9}}{9,00}\times 10^{-6}(2,00 + 4,00) = 5000\ \mathrm{N/C}\\
E_y &= \frac{\sqrt{2,75}}{3,00}\times\frac{8,99\times 10^{9}}{9,00}\times 10^{-6}(4,00 - 2,00) = 1105\ \mathrm{N/C}
\end{align*}

\punts{0,2}
\[ \left.\begin{aligned}
&\left|\vec{E}_\mathit{total}\right| = \sqrt{5000^2 + 1105^2} = 5120\ \mathrm{N/C}\\
&\phi = \mathrm{arc\,tg}\,\frac{1105}{5000} = 12,5°\ \text{respecte a l'eix X}
\end{aligned}\right\}\ \vec{E}_\mathit{total} = \left(5000\,\vec{\imath} + 1105\,\vec{\jmath}\right)\ \mathrm{N/C} \]

\punts{0,1} $V_T = V_1 + V_2;\ V = k\dfrac{Q}{r}$

\punts{0,2} $V_T = \dfrac{8,99\times 10^{9}}{3,00}\times 10^{-6}(2,00 - 4,00) = -6000\ \mathrm{V}$
\end{solapartat}

\begin{solapartat}
\punts{0,4} Les dues càrregues han de tenir el mateix signe ja sigui + o $-$\\
D'aquesta manera, els camps $\boldsymbol{E_1}$ i $\boldsymbol{E_2}$ en el punt que indica el problema tindran la mateixa direcció i sentits oposats, podent-se anular.

\punts{0,4}
\begin{align*}
E_1 &= E_2\\
k\frac{Q_1}{r_1^2} &= k\frac{Q_2}{(x - r_1)^2}
\end{align*}

\punts{0,2} $\dfrac{Q_2}{Q_1} = \dfrac{(5 - 1)^2}{1^2} = 16 \Rightarrow Q_2 = 16Q_1$
\end{solapartat}
